\section{Introduction}
\label{sec:introduction}

Topology optimization (TO) provides a systematic numerical design tool for high-performance lightweight structures by optimizing the spatial distribution of material within a prescribed design domain. Since the continuum material distribution approach was introduced by Bends{\o}e~\cite{Bendsoe1989}, density-based topology optimization, particularly represented by the Solid Isotropic Material with Penalization (SIMP) method, has been widely adopted due to its conceptual clarity, algorithmic maturity, and broad engineering applicability~\cite{BendsoeSigmund2003}. In classical implementations, the linear elastic state equations are predominantly discretized using standard conforming displacement-based Lagrange finite element methods (LFEM), where stress tensors are subsequently recovered within element interiors via numerical gradients of the discrete displacement fields. While such a displacement-driven paradigm is computationally efficient and sufficient for conventional compliance minimization problems, it encounters severe discretization errors and numerical bottlenecks in design problems governed by accurate local stress evaluation or nearly incompressible material behavior.

In local stress-constrained problems, constraints are imposed element-wise, whose quantity scales with the mesh size, and their values are strictly dictated by the accuracy of the discrete stress field. In standard displacement elements, differentiating displacement shape functions reduces the stress approximation accuracy by one order (e.g., linear elements yield only piecewise constant stresses), accompanied by non-physical discontinuous jumps in both normal and tangential stresses across element interfaces. If these post-processed discrete stresses are directly substituted into local constraints, the mutual coupling among discretization errors, smoothing strategies, and the optimization model readily induces local stress oscillations or traps the structural evolution in over-optimized designs~\cite{Duysinx1998,Le2010,Paris2010,daSilva2019}. In nearly incompressible material regimes, as Poisson's ratio approaches the incompressible limit ($\nu_0 \to 0.5$), the bulk modulus of plane-strain or three-dimensional continua tends to infinity relative to the shear modulus, imposing a stringent divergence-free deformation constraint ($\operatorname{div}\boldsymbol{u} \approx 0$). Owing to the lack of sufficient deformation degrees of freedom in the discrete displacement space, standard low-order displacement elements suffer from severe volumetric locking~\cite{Puso2006}, which artificially exaggerates structural stiffness and produces pathological thin members and spurious hinges during topological evolution~\cite{SigmundClausen1998,Kumar2017}.

Mixed finite element methods offer an intrinsic variational pathway to overcome the aforementioned mechanical challenges~\cite{BoffiBrezziFortin2013}. Based on the Hellinger--Reissner two-field variational principle, the symmetric stress $\boldsymbol{\sigma}$ and displacement $\boldsymbol{u}$ are treated simultaneously as independent unknowns. The stress trial functions belong directly to the symmetric tensor $H(\operatorname{div})$ space, naturally guaranteeing strict continuity of normal tractions across inter-element boundaries, while the divergence equilibrium equation is satisfied exactly in the weak form. Because the compliance bilinear form remains uniformly bounded in the incompressible limit, mixed elements satisfying the discrete inf-sup (Babu\v{s}ka--Brezzi) condition fundamentally circumvent volumetric locking~\cite{Stenberg1986,BrezziAdams2003}. The truly-mixed topology optimization studies by Bruggi and Venini~\cite{BruggiVenini2007,BruggiVenini2008} and Bruggi~\cite{Bruggi2016} demonstrated that directly discretizing stress and displacement provides an effective means for nearly incompressible structural design and stress-constrained optimization. However, most existing mixed finite element topology optimization literature is restricted to lowest-order discretizations, structured quadrilateral meshes, or specific hybridized formulations. For unstructured simplicial meshes (triangles and tetrahedra), constructing strictly symmetric and conforming $H(\operatorname{div})$ spaces has long posed substantial mathematical difficulties, often necessitating additional multipliers to relax the symmetry into weakly symmetric formulations~\cite{ArnoldWinther2002,AdamsCockburn2005,Arnold2007}.

The Hu--Zhang mixed finite element family, introduced by Hu and Zhang~\cite{HuZhang2014} and Hu~\cite{Hu2015}, represents a landmark breakthrough in the construction of symmetric stress finite elements. By directly constructing full polynomial spaces $\mathbb{P}_k(\mathbb{S})$ on simplices, this method explicitly decouples normal trace and tangential components via a geometric entity decomposition, ensuring that inter-element normal tractions are uniquely governed by interface degrees of freedom, thereby establishing a family of strongly symmetric conforming spaces for arbitrary polynomial degrees. Building upon this foundation, a suite of essential theories has progressively emerged to facilitate its practical computation: Chen et al.~\cite{Chen2017} addressed the high polynomial degree requirement of the native formulation ($k \ge d+1$) by devising a symmetric matrix jump stabilization mechanism for lower orders ($k \le d$); Hu and Ma~\cite{HuMa2021} developed a local stress relaxation theory at complex boundary corners, eliminating over-constraint at geometric singularities; Chen et al.~\cite{Chen2018a,Chen2018b} investigated fast auxiliary space preconditioners and residual-based a posteriori error estimators; and Chen et al.~\cite{Chen2024} subsequently formulated geometric decompositions and efficient algorithms for high-order simplicial face and edge elements. In light of the growing trend toward incorporating high-order finite elements into topology optimization~\cite{ZegardPaulino2016}, how to seamlessly embed such arbitrary-order symmetric stress approximations on simplicial meshes into density-based topology optimization—overcoming the design dependence introduced by non-homogeneous traction lifting, the corner stress relaxation amidst evolving structural boundaries, and the consistent adjoint sensitivity analysis driven by independent stresses—still lacks a systematic variational formulation and computational framework.

To address the aforementioned theoretical and algorithmic bottlenecks, this paper establishes an arbitrary-order Hu--Zhang mixed finite element framework for density-based topology optimization on simplicial meshes, built upon the open-source finite element library FEALPy~\cite{Zheng2026FEALPy} and the modular topology optimization platform SOPTX~\cite{He2026SOPTX}. The principal contributions of this work are summarized as follows:
\begin{enumerate}
    \item \textbf{Variational formulation and consistent adjoint analysis for simplicial mixed topology optimization}: Starting from the Hellinger--Reissner variational principle with non-homogeneous Neumann boundary conditions, an extrinsic traction lifting field is employed to handle boundary traction data, systematically revealing the design dependence of the load vector induced by the lifting field. Based on the reconstructed total stress, we derive a consistent adjoint sensitivity analysis framework for both complementary energy objectives and augmented Lagrangian local stress constraints.
    \item \textbf{Consistent degree-of-freedom management and operator reuse for optimization iterations}: By casting the intricate high-order tensor degrees of freedom (DOFs) into a unified management of normal trace continuity, and integrating a global reference coordinate frame with local corner stress relaxation at geometric singularities, we construct an efficient computational implementation wherein the discrete divergence and stabilization operators are fully reused both structurally and numerically throughout the optimization process, requiring dynamic updates exclusively for the compliance blocks.
    \item \textbf{Cross-order numerical evaluation and mechanical applicability boundaries}: We design a comprehensive benchmark suite spanning low- and high-order manufactured solution convergence rates, cross-order compliance consistency between displacement and mixed formulations, volumetric locking alleviation in nearly incompressible regimes, and denominator-free local stress-constrained optimization. This objectively delineates the practical applicability boundaries of independent symmetric stress approximations relative to standard displacement elements in terms of topological evolution, stress fidelity, and computational cost.
\end{enumerate}

The remainder of this paper is organized as follows. Section~\ref{sec:preliminaries} introduces the governing equations of linear elasticity, the Hellinger--Reissner mixed variational formulation, and the traction lifting field at the continuum level. Section~\ref{sec:huzhang_elements} elaborates on the tensor geometric decomposition on simplices, the construction of arbitrary-order conforming spaces, jump stabilization for lower polynomial degrees, local corner stress relaxation, and the assembly of discrete saddle-point systems. Section~\ref{sec:topopt_formulation} establishes the two-parameter material interpolation, the complementary energy compliance minimization model, singularity-free apparent stress constraints, and consistent adjoint sensitivities based on the total stress. Section~\ref{sec:numerical_results} presents convergence rate verifications with manufactured solutions alongside three representative benchmark problems in topology optimization. Finally, Section~\ref{sec:conclusions} summarizes the conclusions and delineates the operational boundaries of the proposed framework.
